\section{Introduzione}

\subsection{Executive Summary}
Il presente documento costituisce il deliverable D1 del Progetto di Ingegneria del Software 2026/2027
e riporta lo studio di fattibilità dell'app Android a supporto delle attività del Laboratorio in Campo
degli studenti del secondo anno di Scienze Ambientali, sviluppata dal gruppo \emph{I paguri}.
Il documento descrive il coinvolgimento degli stakeholder, svolto tramite un questionario a studenti e docenti (7 risposte), e i primi requisiti che ne derivano. Analizza poi tre applicazioni simili (QField, Rock Identifier, KoboCollect) e i due mock-up dell'hackathon di settembre, e individua gli strumenti hardware e software preliminari per lo sviluppo. Dal questionario emerge che il problema principale è il tempo necessario per trascrivere i dati raccolti sul campo, spesso in zone senza rete, e che le funzioni più richieste sono foto e note vocali collegate alla posizione, la ricerca delle specie e l'uso offline.

\subsection{Struttura del documento}
La Sezione~2 descrive il contesto e gli obiettivi. La Sezione~3 riporta tecniche, modalità e tempistica
per il coinvolgimento degli stakeholder e i primi risultati sui requisiti. La Sezione~4 analizza le
applicazioni simili e i mock-up dell'hackathon di settembre. La Sezione~5 individua gli strumenti
(hardware e software) per lo sviluppo. La Sezione~6 riporta le conclusioni.
